\documentclass[aspectratio=169,dvipdfmx,11pt,notheorems]{beamer}
%%%% 和文用 %%%%%
\usepackage{bxdpx-beamer}
\usepackage{pxjahyper}
\usepackage{minijs}%和文用
\renewcommand{\kanjifamilydefault}{\gtdefault}%和文用

%%%% スライドの見た目 %%%%%
\usetheme{Madrid}
\usefonttheme{professionalfonts}
\setbeamertemplate{frametitle}[default][center]
\setbeamertemplate{navigation symbols}{}
\setbeamercovered{transparent}%好みに応じてどうぞ）
\setbeamertemplate{blocks}[rounded]
\useinnertheme{circles}
\setbeamertemplate{footline}[page number]
\setbeamerfont{footline}{size=\normalsize,series=\bfseries}
\setbeamercolor{footline}{fg=black,bg=black}
%%%%

%%%% 定義環境 %%%%%
\usepackage{amsmath,amssymb}
\usepackage{amsthm}
\theoremstyle{definition}
\newtheorem{theorem}{定理}
\newtheorem{definition}{定義}
\newtheorem{proposition}{命題}
\newtheorem{lemma}{補題}
\newtheorem{corollary}{系}
\newtheorem{conjecture}{予想}
\newtheorem*{remark}{Remark}
\renewcommand{\proofname}{}
%%%%%%%%%

%%%%% フォント基本設定 %%%%%
\usepackage[T1]{fontenc}%8bit フォント
\usepackage{textcomp}%欧文フォントの追加
\usepackage[utf8]{inputenc}%文字コードをUTF-8
\usepackage[deluxe]{otf}%otfパッケージ
\usepackage{lxfonts}%数式・英文ローマン体を Lxfont にする
\usepackage{bm}%数式太字
%%%%%%%%%%

%%%%% PythonTeX %%%%%
\usepackage[makestderr]{pythontex}
\restartpythontexsession{\thesection}

\usepackage[dvipdfmx]{media9}
\usepackage{ulem}
 
\title{Let's implement useless Python objects}
\author[Hayao]{Hayao Suzuki}
\institute[PythonAsia 2026]{PythonAsia 2026 at DE LA SALLE UNIVERSITY, MALATE, MANILA}
%\titlegraphic{\includegraphics[scale=0.2]{pyconlogo.png}}
\date{March 21, 2026}

\begin{document}

\begin{frame}[plain]\frametitle{}
\titlepage %表紙
\includemedia[
  addresource=voice/slide01.wav,
  flashvars={
    source=voice/slide01.wav
   &autoPlay=true
  },
  transparent
]{\color{blue}{\tiny\fbox{Play}}}{APlayer.swf}
\end{frame}

\begin{frame}\frametitle{Share it}

\begin{block}{GitHub}
\begin{itemize}
\item \url{https://github.com/HayaoSuzuki/pythonasia2026}
\end{itemize}
\end{block}

\begin{block}{Hashtag}
\begin{itemize}
\item \#PythonAsia2026
\end{itemize}
\end{block}
\includemedia[
  addresource=voice/slide02.wav,
  flashvars={
    source=voice/slide02.wav
   &autoPlay=true
  },
  transparent
]{\color{blue}{\tiny\fbox{Play}}}{APlayer.swf}
\end{frame}

\section{Self Introduction}

\begin{frame}\frametitle{Who am I ?}

\begin{block}{お前誰よ}
\begin{description}
\item[Name] Hayao Suzuki（鈴木　駿）
\item[\xout{Twitter} X] \href{https://twitter.com/CardinalXaro}{@CardinalXaro}
\item[Work] Software Engineer at TOKYO GAS CO.,LTD
\end{description}
\end{block}

\begin{block}{About Tokyo Gas}
Tokyo Gas is the largest natural gas utility in Japan
\begin{itemize}
\item Production and sale of city gas, LNG sales, etc.
\item Gas Pipeline Service Business、Supply of city gas
\item Production, supply and sale of electricity
\end{itemize}
\end{block}
\includemedia[
  addresource=voice/slide03.wav,
  flashvars={
    source=voice/slide03.wav
   &autoPlay=true
  },
  transparent
]{\color{blue}{\tiny\fbox{Play}}}{APlayer.swf}


\end{frame}

\begin{frame}\frametitle{Who am I ?}

\begin{block}{Presentations of PyCon}
\begin{itemize}
\item \structure{明日からgraphlib、みんなで使おう}(PyCon JP 2025)
\item \structure{Let's implement useless Python objects}(PyCon APAC 2023)
\item \structure{組み込み関数powの知られざる進化}(PyCon JP 2021)
\item \structure{インメモリーストリーム活用術}(PyCon JP 2020)
\item \structure{君はcmathを知っているか}(PyCon mini Shizuoka 2020)
\item \structure{Pythonと楽しむ初等整数論}(PyCon mini Hiroshima 2019)
\item \structure{SymPyによる数式処理}(PyCon JP 2018)
\end{itemize}
\end{block}
Listed at  \url{https://xaro.hatenablog.jp/}

\includemedia[
  addresource=voice/slide05.wav,
  flashvars={
    source=voice/slide05.wav
   &autoPlay=true
  },
  transparent
]{\color{blue}{\tiny\fbox{Play}}}{APlayer.swf}
\end{frame}

\section{Useless!}

\begin{frame}\frametitle{Today's Theme}

\begin{center}
\huge{Let's implement \structure{useless} Python objects}
\end{center}
\includemedia[
  addresource=voice/slide06.wav,
  flashvars={
    source=voice/slide06.wav
   &autoPlay=true
  },
  transparent
]{\color{blue}{\tiny\fbox{Play}}}{APlayer.swf}
\end{frame}

\begin{frame}\frametitle{What is it mean useless?}

\begin{block}{From  Longman Dictionary of Contemporary English}
\begin{enumerate}
\item not useful or effective in any way
\item (informal) unable or unwilling to do anything properly
\end{enumerate}
\end{block}
\includemedia[
  addresource=voice/slide07.wav,
  flashvars={
    source=voice/slide07.wav
   &autoPlay=true
  },
  transparent
]{\color{blue}{\tiny\fbox{Play}}}{APlayer.swf}
\end{frame}

\begin{frame}\frametitle{Is the useless object really useless?}

\begin{block}{From Zhuangzi "The Human World"（荘子 人間世篇）}
人皆知有用之用 而莫知無用之用也
\end{block}
Everyone knows the usefulness of the useful, but no one knows the usefulness of the useless.

\includemedia[
  addresource=voice/slide08.wav,
  flashvars={
    source=voice/slide08.wav
   &autoPlay=true
  },
  transparent
]{\color{blue}{\tiny\fbox{Play}}}{APlayer.swf}
\end{frame}

\begin{frame}\frametitle{Today's Theme}

\begin{block}{Let's implement useless Python objects}
The useless objects are useless, but \structure{how to make a useless object} is very useful.
\end{block}
\includemedia[
  addresource=voice/slide09.wav,
  flashvars={
    source=voice/slide09.wav
   &autoPlay=true
  },
  transparent
]{\color{blue}{\tiny\fbox{Play}}}{APlayer.swf}
\end{frame}

\section{Basic useless Python objects}

\begin{frame}[fragile]\frametitle{What is a useless Python object?}

\begin{exampleblock}{Example: \texttt{LiarContainer}}
\begin{pygments}{python}
>>> c = LiarContainer(["spam", "egg", "bacon"])
>>> "spam" in c
False
>>> "tomato" in c
True
\end{pygments}
\end{exampleblock}

\end{frame}
\begin{frame}[fragile]\frametitle{What is a useless Python object?}

\begin{exampleblock}{Example: \texttt{FibonacciSized}}
\begin{pygments}{python}
>>> s = FibonacciSized(range(50))
>>> len(s)
12586269025
\end{pygments}
\end{exampleblock}

\end{frame}

\begin{frame}[fragile]\frametitle{What is a useless Python object?}

\begin{exampleblock}{Example: \texttt{ShuffledIterable}}
\begin{pygments}{python}
>>> it = ShuffledIterable([1, 2, 3, 4, 5])
>>> for _ in range(3):
...     for v in it:
...         print(v, end=" ")
...     print()
...
5 3 4 2 1 
4 1 2 3 5 
2 5 3 1 4
\end{pygments}
\end{exampleblock}

\end{frame}

\begin{frame}\frametitle{What is a useless Python object?}

\begin{block}{Definition of a useless Python object in this talk}
A useless Python object behave \structure{Pythonic}, but does not work as expected.
\end{block}

\end{frame}

\begin{frame}[fragile]\frametitle{Data Structures and Operations}

\begin{block}{Basic Data Structures of Python}
\begin{description}
\item[List] \pyv{[1, 4, 9, 16, 25, 36]}
\item[Tuple] \pyv{("pen", "pineapple", "apple", "pen")}
\item[Dictionary] \pyv{{"Answer": 42}}
\item[Set] \pyv{{41, 43, 47, 53, 57, 59}}
\end{description}
\end{block}

\begin{block}{Common Operations of Data Structure}
\begin{description}
\item[\texttt{len()}] Length of object
\item[\texttt{in}] Membership test
\item[\texttt{for}] Iteration
\end{description}
\end{block}

\end{frame}

%\begin{frame}[fragile]\frametitle{\texttt{Useless} Abstract Base Class}
%
%\begin{exampleblock}{Example: \texttt{Useless} ABC}
%\begin{pygments}{python}
%class Useless(abc.ABC):
%    def __init__(self, data: Optional[Iterable] = None):
%        if data is not None:
%            self._data = [v for v in data]
%        else:
%            self._data = []
%\end{pygments}
%\end{exampleblock}
%\texttt{Useless} abstract base is useful, contrary to its name.
%\end{frame}

\begin{frame}[fragile]\frametitle{\texttt{in} and \texttt{Container}}

\begin{block}{\texttt{object.\_\_contains\_\_()}}
Called to implement membership test operators.
\end{block}

\begin{exampleblock}{Example: \texttt{LiarContainer}}
\begin{pygments}{python}
class LiarContainer(Container):
    def __contains__(self, item) -> bool:
        return item not in self._data
\end{pygments}
\end{exampleblock}

\end{frame}

\begin{frame}[fragile]\frametitle{\texttt{len()} and \texttt{Sized}}

\begin{block}{\texttt{object.\_\_len\_\_()}}
Called to implement the built-in function \texttt{len()}.
\end{block}

\begin{exampleblock}{Example: \texttt{FibonacciSized}}
\begin{pygments}{python}
class FibonacciSized(Sized):
    PHI: Final[float] = (1 + math.sqrt(5)) / 2
    def __len__(self) -> int:
        return round(
            (1 / math.sqrt(5))
            * pow(self.PHI, len(self._data))
        )
\end{pygments}
\end{exampleblock}

\end{frame}

\begin{frame}[fragile]\frametitle{\texttt{for} and \texttt{Iterable}}

\begin{block}{\texttt{object.\_\_iter\_\_()}}
Called when an iterator is required for a container. 
\end{block}

\begin{exampleblock}{Example: \texttt{ShuffledIterable}}
\begin{pygments}{python}
class ShuffledIterable(Iterable):
    def __iter__(self) -> Iterator:
        return iter(
            random.sample(self._data, k=len(self._data))
        )
\end{pygments}
\end{exampleblock}

\end{frame}

\begin{frame}[fragile]\frametitle{\texttt{reversed()} and \texttt{Reversible}}

\begin{block}{\texttt{object.\_\_reversed\_\_()}}
Called to implement the built-in function \texttt{reversed()}.
\end{block}

\begin{exampleblock}{Example: \texttt{ReversedReversible}}
\begin{pygments}{python}
class ReversedReversible(Reversible):
    def __iter__(self) -> Iterator:
        return iter(reversed(self._data))

    def __reversed__(self) -> Iterator:
        return iter(self._data)
\end{pygments}
\end{exampleblock}

\end{frame}

\begin{frame}\frametitle{Object Protocols}

\begin{block}{How to implement Pythonic Python objects}
We need to understand \structure{object protocols}.
\end{block}
Ref: \url{https://docs.python.org/3/reference/datamodel.html}

\end{frame}

\begin{frame}\frametitle{\texttt{collections.abc}}

\begin{block}{\texttt{collections.abc}}
This module provides \structure{abstract base classes} that can be used to test whether a class provides \structure{a particular interface}.
\end{block}
From \url{https://docs.python.org/3/library/collections.abc.html}

\end{frame}

\begin{frame}\frametitle{\texttt{collections.abc}}

\begin{exampleblock}{\texttt{collections.abc} and Interface}
\begin{table}[h]
\centering
\begin{tabular}{lcr}
\hline
ABC  & Interface \\
\hline
\texttt{Sized}  & \texttt{\_\_len\_\_()} \\
\texttt{Container}  & \texttt{\_\_contains\_\_()} \\
\texttt{Iterable}  &  \texttt{\_\_iter\_\_()} \\
\texttt{Reversible}  &  \texttt{\_\_reversed\_\_()} \\
\texttt{Collection}  & \texttt{Sized}, \texttt{Container}, \texttt{Iterable}  \\
\hline
\end{tabular}
\end{table}
\end{exampleblock}

\end{frame}

\begin{frame}[fragile]\frametitle{\texttt{Collection}}

\begin{block}{\texttt{collections.abc.Collection}}
\texttt{Sized} and \texttt{Container} and \texttt{Iterable}
\end{block}

\begin{exampleblock}{Example: \texttt{Collection}}
\begin{pygments}{python}
class UselessCollection(
    FibonacciSized, LiarContainer, ShuffledIterable
):
    pass
\end{pygments}
\end{exampleblock}

\end{frame}

\begin{frame}\frametitle{\texttt{collections.abc}}

\section{Practical useless Python objects}

\begin{exampleblock}{\texttt{collections.abc} and Built-in Objects}
\begin{table}[h]
\centering
\begin{tabular}{lcr}
\hline
ABC  & built-in objects \\
\hline
\texttt{Sequence}  & \texttt{tuple} \\
\texttt{MutableSequence}  & \texttt{list} \\
\texttt{MutableSet}  &  \texttt{set} \\
\texttt{MutableMapping}  & \texttt{dict}  \\
\hline
\end{tabular}
\end{table}
\end{exampleblock}

\end{frame}

\begin{frame}[fragile,shrink=10]\frametitle{\texttt{Sequence}}

\begin{exampleblock}{Example: \texttt{ModularSequence}}
\begin{pygments}{python}
class ModularSequence(Sequence):
    def __getitem__(self, key):
        match key:
            case int():
                return self._data[key % len(self._data)]
            case slice():
                n = len(self._data)
                s = slice(
                    key.start % n if key.start is not None else None,
                    key.stop % n if key.stop is not None else None,
                    key.step,
                )
                return self._data[s]
            case _:
                raise TypeError
\end{pygments}
\end{exampleblock}

\end{frame}

\begin{frame}[fragile]\frametitle{\texttt{Sequence}}

\begin{exampleblock}{Example: \texttt{ModularSequence}}
\begin{pygments}{python}
>>> seq = ModularSequence(range(20))
>>> print(seq[21:44])
[1, 2, 3]
>>> print(seq[65543])
3
>>> seq.count(13) # It does not stop.
\end{pygments}
\end{exampleblock}

\end{frame}

\begin{frame}[fragile]\frametitle{\texttt{Sequence}}

\begin{exampleblock}{Example: \texttt{CompetitionSequence}}
\begin{pygments}{python}
class CompetitionSequence(Sequence):
    def __getitem__(self, index):
        return self._data[index]

    def __iter__(self):
        return iter(reversed(self._data))

    def __len__(self):
        return len(self._data)
\end{pygments}
\end{exampleblock}

\end{frame}

\begin{frame}[fragile]\frametitle{\texttt{Sequence}}

\begin{exampleblock}{Example: \texttt{CompetitionSequence}}
\begin{pygments}{python}
>>> s = CompetitionSequence("abcdefg")
>>> s[0]
'a'
>>> list(s)
['g', 'f', 'e', 'd', 'c', 'b', 'a']
>>> len(s)
7
\end{pygments}
\end{exampleblock}

\end{frame}

\begin{frame}[fragile]\frametitle{\texttt{Mapping}}

\begin{exampleblock}{Example: \texttt{MisprintedDictionary}}
\begin{pygments}{python}
class MisprintedDictionary(Mapping):
    def __init__(self, _dict: dict):
        keys = list(_dict.keys())
        values = list(_dict.values())
        rotated_values = chain(
            islice(values, 1, None),
            islice(values, 1),
        )
        self._data = dict(
            zip(keys, rotated_values)
        )
\end{pygments}
\end{exampleblock}

\end{frame}

\begin{frame}[fragile]\frametitle{\texttt{Mapping}}

\begin{exampleblock}{Example: \texttt{MisprintedDictionary}}
\begin{pygments}{python}
>>> d = MisprintedDictionary({"a": 1, "b": 2, "c": 3})
>>> d["a"], d["b"], d["c"]
(2, 3, 1)
>>> list(d.keys())
['a', 'b', 'c']
>>> len(d)
2
\end{pygments}
\end{exampleblock}

\end{frame}

\begin{frame}[fragile]\frametitle{\texttt{Set}}

\begin{exampleblock}{Example: \texttt{CrowdSet}}
\begin{pygments}{python}
@functools.total_ordering
class CrowdSet(Set):
    def __init__(self, data=None):
        if data is not None:
            self._data = set(data)
        else:
            self._data = set()

    def __lt__(self, other):
        return self._data >= other._data
\end{pygments}
\end{exampleblock}

\end{frame}

\begin{frame}[fragile]\frametitle{\texttt{Set}}

\begin{exampleblock}{Example: \texttt{CrowdSet}}
\begin{pygments}{python}
class CrowdSet(Set):
    def __len__(self):
        return len(self._data) ** 2

    def __and__(self, other):
        return CrowdSet(self._data | other._data)

    def __or__(self, other):
        return CrowdSet(self._data & other._data)
\end{pygments}
\end{exampleblock}

\end{frame}

\begin{frame}[fragile]\frametitle{\texttt{Set}}

\begin{exampleblock}{Example: \texttt{CrowdSet}}
\begin{pygments}{python}
>>> s = CrowdSet(("egg", "bacon", "spam"))
>>> t = CrowdSet(("egg", "egg", "spam", "spam"))
>>> s > t
True
>>> len(s)
9
>>> sorted(s & t)
['bacon', 'egg', 'spam']
>>> sorted(s | t)
['egg', 'spam']
\end{pygments}
\end{exampleblock}

\end{frame}

\section{Conclusion}

\begin{frame}\frametitle{Conclusion}

\begin{block}{Let's implement useless Python objects}
\begin{itemize}
\item Useless Python objects are useful
\item \texttt{collections.abc} module is very useful
\item Once you understand object protocol, you can do anything
\end{itemize}
\end{block}
\end{frame}



\end{document}